% sections/D_classical.tex: Lane m2 (classical machine learning on engineered features), appendix.
% Numbers only through \R / \Rci; figures through \genfig, tables through \gentab (report/README.md).
% Sources: analysis/reports/stage2/m2_SUMMARY.md, figures/m2/captions.json, tables/m2/*, results/m2_numbers.json.
\section{Classical machine learning on engineered features}\label{app:m2}

Lane m2 asks how far classical (non-deep) models get on a rich set of hand-crafted per-trial features, and
which data treatment and which features matter. Every result is leave-one-person-out over the
\R{describe/n_persons} persons and \R{describe/primary_trials} primary trials. Every hyperparameter is tuned
inside the outer fold, on \R{m2/inner_folds} inner person folds of the training persons. A paired difference
compares two pipelines on the same trials, with a paired person-bootstrap CI. Its p-value is an exact sign-flip
test over persons, Holm-adjusted within the comparison family named in the text (level \R{m2/alpha}). Two
pipelines whose \pright{} differ by less than \R{m2/equiv_dp} on every trial are called identical (solver
precision); they are not tested. The inputs are the oracle approach windows of the cache
(Appendix~\ref{app:prep}); production sees detected windows instead (Appendix~\ref{app:prod}).
Abbreviations used in this appendix: logistic regression (LR), support vector machine (SVM), radial basis
function (RBF), k nearest neighbours (k-NN), LightGBM gradient-boosted trees (LGBM), calibration-in-the-large
(CITL).

% ------------------------------------------------------------------------------------------------------------
\subsection{Features and their mirror parity}\label{app:m2:features}

\paragraph{Representation.} One trial is described by \R{m2/lr_full_n_features} features computed from the
17 channels of the cache. Nine are in the device frame: the three axes of the linear acceleration, of the
rotation rate and of the unit up vector. Eight are in the gravity frame: magnitudes, vertical and horizontal
components, the yaw rate about gravity $\vomega\cdot\ghat$, its running integral (the cumulative yaw) and the
tilt from the start. The approach window is time-normalised to 64 samples from onset to end. Seven feature
families describe it (Table~\ref{tab:m2_feature_families}): means in $k = 8$ equal time bins (\code{bins}),
summary statistics (\code{stats}), log band powers and the lowest Fourier coefficients (\code{fft}), wavelet
sub-band energies (\code{wavelet}), correlations of every channel pair within a frame group (\code{xcorr}),
cumulative integrals at the quarters of the window (\code{cumint}) and the log approach duration
(\code{global}).

\paragraph{Mirror parity.} Every feature is even (unchanged), odd (sign flip) or one of a swapped pair under the
left-right mirror $m(\cdot)$ of Eqs.~\eqref{eq:mirror} and~\eqref{eq:parity}. The feature registry gives the
exact mirror map of every feature (unit-tested): the features of a mirrored trial are a signed permutation of
the original ones. Of the default set, \R{m2/reg_odd} features are odd or swapped and \R{m2/reg_even} are even.
The \code{wavelet} and \code{global} families are entirely even (\R{m2/reg_wavelet_odd} and
\R{m2/reg_global_odd} odd features). This parity matters for the design: with mirror augmentation each even
feature value enters the training set once with each label, so a linear model gives every even feature (and the
intercept) zero weight. The augmented LR therefore uses only the mirror-sensitive features.

\gentab{m2/m2_feature_families}

% ------------------------------------------------------------------------------------------------------------
\subsection{Models, tuning and the headline estimate}\label{app:m2:models}

\paragraph{Options.} The lane compares \R{m2/n_models_compared} pipelines on the same features, all
mirror-augmented: the pre-specified
default, an L2-penalised LR; L1 and elastic-net LR; a linear SVM, an RBF SVM and a k-NN classifier, each with a
Platt map \cite{platt1999probabilistic}; a random forest and extra trees with one fixed configuration
(\R{m2/forest_n_estimators} trees, $\sqrt{p}$ of the $p$ features per split, leaves of at least
\R{m2/forest_min_leaf} rows; not tuned); HistGradientBoosting; LGBM; and a soft vote and a stack of the L2 LR,
the RBF SVM and LGBM. The models are those of scikit-learn \cite{pedregosa2011sklearn} and LightGBM. The L2
and L1 LR and the linear SVM are tuned on a grid of their penalty. The other tuned models use a small seeded
Bayesian search (Optuna) per outer fold: \R{m2/trials_svm_rbf} trials for the RBF SVM, \R{m2/trials_knn} for
the k-NN, \R{m2/trials_lgbm} for LGBM, \R{m2/trials_lr_en} for the elastic net and \R{m2/trials_hgb} for
HistGradientBoosting.

\paragraph{Three estimates with three roles.} The \emph{nested selection} (\code{m2_select}) chooses, inside
every outer fold, among the L2 LR, the RBF SVM, LGBM and their soft vote by inner cross-validated log loss,
then refits the winner on all training persons. It is one procedure, so its leave-one-person-out estimate is
selection-aware: it is the lane's headline. The \emph{default} (\code{m2_lr_full}) was fixed before any outer
result. The best of the \R{m2/n_models_compared} pipelines by outer log loss (the stack) and the best tree
pipeline (LGBM) are \emph{exploratory}: each is the best of \R{m2/ens_stack_configs} (stack) and
\R{m2/lgbm_configs} (LGBM) configurations compared on the outer folds. The selection chose the soft vote in \R{m2/sel_choice_ens_soft} and LGBM in
\R{m2/sel_choice_lgbm} of the \R{m2/sel_n_folds} outer folds. The median inner log-loss margin over the
runner-up is \R{m2/sel_margin_median}, so the candidates are close. Every candidate's inner predictions are
out of sample; the SVM's Platt map is cross-fitted across the inner folds, which corrects an in-sample map
found by the audit.

\paragraph{Results.} Figure~\ref{fig:M2_01_models} and Table~\ref{tab:m2_models} show every pipeline. The
nested selection reaches an accuracy of \Rci{m2/select_accuracy}, an AUC of \Rci{m2/select_auc}, a log loss of
\Rci{m2/select_log_loss} and a Brier score of \Rci{m2/select_brier}. The default LR reaches
\Rci{m2/lr_full_accuracy} with a log loss of \Rci{m2/lr_full_log_loss}. The reference floor, a mirror-augmented
LR on the cumulative yaw at 8 instants, reaches \Rci{m2/ref_accuracy}. The selection beats it by
\Rcib{m2/select_ref_d_acc}~points (Holm $p = \R{m2/select_ref_p_holm}$) and the default by
\Rcib{m2/lr_full_ref_d_acc}~points (Holm $p = \R{m2/lr_full_ref_p_holm}$). Both are adjusted over the
\R{m2/fam_vs_reference_size} comparisons of m2 pipelines with the reference floor. Every pipeline of this family
beats the floor after the adjustment. The smallest gain is that of the L1 LR (\Rcib{m2/lr_l1_ref_d_acc}~points,
Holm $p = \R{m2/lr_l1_ref_p_holm}$), the largest that of the stack (\Rcib{m2/ens_stack_ref_d_acc}~points, Holm
$p = \R{m2/ens_stack_ref_p_holm}$).

\paragraph{Against the default.} The family is the \R{m2/fam_model_vs_default_size} comparisons of the other
pipelines and the selection with the default. No pipeline differs from the default in accuracy after the Holm
adjustment. The stack, the best pipeline by outer log loss, gains \Rcib{m2/ens_stack_d_acc}~points of accuracy
(Holm $p = \R{m2/ens_stack_p_acc_holm}$) and changes the log loss by \Rcib{m2/ens_stack_d_ll} (Holm
$p = \R{m2/ens_stack_p_ll_holm}$): not significant. The soft vote is the only pipeline with a lower log loss
after the adjustment (\Rcib{m2/ens_soft_d_ll}, Holm $p = \R{m2/ens_soft_p_ll_holm}$), and no pipeline has a
higher one. The nested selection changes the log loss by \Rcib{m2/select_d_ll} (Holm
$p = \R{m2/select_p_ll_holm}$). With \R{describe/n_persons} persons, the spread of a few points between
pipelines is not resolved.

\paragraph{Natural use and abstentions.} The free-choice trials are all right-hand trials so far
(\R{m2/nat_n_right} right, \R{m2/nat_n_left} left). Their score is therefore right-hand recall, not an accuracy
over both hands, and a constant `right' answer would score 100\,\%. The nested selection recognises
\R{m2/select_nat_correct} of \R{m2/select_nat_n} free-choice trials (right-hand recall \R{m2/select_nat_acc}
[\Rb{m2/select_nat_lo}, \Rb{m2/select_nat_hi}]), against a right-hand recall of \R{m2/select_right_recall} on the
primary trials. The k-NN and the RBF SVM recognise \R{m2/knn_nat_correct} (Table~\ref{tab:m2_models}). None of
the four reported pipelines abstains: none returns $\pright = 0.5$ exactly on any primary or free-choice trial
(free-choice coverage \R{m2/select_nat_coverage} for the selection, \R{m2/lr_full_nat_coverage} for the default,
\R{m2/ens_stack_nat_coverage} for the stack and \R{m2/lgbm_nat_coverage} for LGBM). Their answered and
forced-choice accuracies therefore equal the accuracies shown. The only abstaining configuration of the lane is
an ablation (Appendix~\ref{app:m2:families}).

\paragraph{Finalists.} Lane m2 hands three finalists to the cross-lane nested selection of lane final
(Table~\ref{tab:m2_finalists}; Appendix~\ref{app:final}): the best linear, tree and ensemble pipelines by
outer log loss. This choice is exploratory. For \code{m2_lr_full} the two records disagree: the finalists file
counts \R{final/nconf_file_m2_lr} configurations compared on the outer folds, the saved result
\R{final/nconf_result_m2_lr}. Lane final uses the finalists file's count, the role in which the model became a
candidate, so in the cross-lane selection it counts as exploratory. Inside lane m2 it remains the pre-specified
default.

\begin{figure}[tbp]
\centering
\genfig{m2/M2_01_models}
\caption{Classical pipelines on the same \R{m2/lr_full_n_features} features, leave-one-person-out over
\R{describe/n_persons} persons, sorted by log loss (best at the top). Blue: the pre-specified default (L2 LR);
orange: the nested selection, the selection-aware headline; dashed: the reference floor; * fixed
hyperparameters. The best pipelines are each the best of \R{m2/ens_stack_configs} compared on the outer folds:
exploratory.}
\label{fig:M2_01_models}\label{fig:app:m2-models}
\end{figure}

\gentab{m2/m2_models}

\gentab{m2/m2_finalists}

\FloatBarrier
% ------------------------------------------------------------------------------------------------------------
\subsection{Data treatment}\label{app:m2:treatment}

\paragraph{Design.} The default treatment was pre-specified: the time-normalised approach window, all 17
channels, the seven families with $k = 8$ bins, a z-score fitted on the training rows, mirror augmentation and
the L2 LR. Each of \R{m2/fam_data_treatment_size} alternatives is paired against it on the same trials and
Holm-adjusted within this family of \R{m2/fam_data_treatment_size} comparisons. The rule: adopt an alternative
only if its held-out log loss is lower after the adjustment. No alternative has a lower log loss after the
adjustment (Figure~\ref{fig:M2_02_data_treatment}; Table~\ref{tab:m2_data_treatment}, upper rows), so the
default stays.
The mirror options are treated in Appendix~\ref{app:m2:mirror}.

\paragraph{Window.} The alternatives are a fixed \R{m2/fixed_window_s}-s window in absolute time (from 0.5~s
before to 2.5~s after the onset) and both windows side by side. Neither differs from the approach window (fixed:
log loss \Rcib{m2/lr_win_fixed_d_ll}, Holm $p = \R{m2/lr_win_fixed_p_ll_holm}$; both:
\Rcib{m2/lr_win_both_d_ll}, Holm $p = \R{m2/lr_win_both_p_ll_holm}$). The fixed window is also not
onset-only. The cache freezes every channel from the app's end cue on, which happens in \R{m2/fixed_n_frozen}
of \R{m2/fixed_n_trials} trials (median \R{m2/fixed_frozen_share} of the post-onset window frozen; Spearman
correlation of the freeze point with the approach duration \R{m2/fixed_rho}). Its gravity reference also uses
the rest after the end. Production has no end cue. Decision: keep the approach window.

\paragraph{Channel frame.} Keeping only the device-frame channels worsens the log loss by
\Rcib{m2/lr_frame_device_d_ll} (Holm $p$~\R{m2/lr_frame_device_p_ll_holm}) and the accuracy by \Rcib{m2/lr_frame_device_d_acc}~points
(Holm $p = \R{m2/lr_frame_device_p_acc_holm}$). Keeping only the gravity-frame channels worsens the log loss by
\Rcib{m2/lr_frame_gravity_d_ll} (Holm $p = \R{m2/lr_frame_gravity_p_ll_holm}$); its accuracy does not differ
(Holm $p = \R{m2/lr_frame_gravity_p_acc_holm}$). Decision: keep both frames. The device-frame features do not
transfer to unseen grips, however (Appendix~\ref{app:m2:importance}).

\paragraph{Time bins.} With $k = 4$ or $k = 16$ bins instead of 8, the log loss does not differ (Holm
$p = \R{m2/lr_bins4_p_ll_holm}$ and \R{m2/lr_bins16_p_ll_holm}). Decision: keep $k = 8$.

\paragraph{Scaling.} Without scaling, the log loss worsens by \Rcib{m2/lr_scale_none_d_ll} (Holm $p$~\R{m2/lr_scale_none_p_ll_holm}).
The unscaled fit does not converge: the feature standard deviations span \R{m2/feature_sd_orders} orders of
magnitude, and with the solver capped at 300 iterations \R{m2/unscaled_inner_nonconverged} of
\R{m2/unscaled_inner_fits} inner fits and \R{m2/unscaled_final_nonconverged} of \R{m2/sel_n_folds} outer-fold
fits stopped at the cap. Robust scaling (median and interquartile range) worsens the log loss by
\Rcib{m2/lr_scale_robust_d_ll} (Holm $p = \R{m2/lr_scale_robust_p_ll_holm}$). Winsorising at 3 standard
deviations and per-trial amplitude normalisation do not differ (Holm $p = \R{m2/lr_winsor3_p_ll_holm}$ and
\R{m2/lr_ampnorm_p_ll_holm}). Decision: keep the z-score. Trees are scale-free; the linear, kernel and k-NN
models need the scaling.

\paragraph{Inherited treatment.} Filtering, resampling and gap handling are the cache's: a 15~Hz zero-phase
low-pass filter, a 60~Hz grid, interpolation of gaps up to 100~ms and exclusion of longer gaps near the
movement (Appendix~\ref{app:prep}). This lane inherits them unchanged. Every feature is checked finite before a
fit.

\begin{figure}[tbp]
\centering
\genfig{m2/M2_02_data_treatment}
\caption{Data-treatment alternatives of the default L2 LR, each paired against the default on the same trials:
accuracy (left) and log loss (right) differences with paired 95\,\% CIs. Orange: Holm-adjusted $p < 0.05$ within
the family of \R{m2/fam_data_treatment_size} comparisons; open marker: identical within solver precision (not a
test). No alternative improves the log loss, so the pre-specified default stays.}
\label{fig:M2_02_data_treatment}\label{fig:app:m2-treatment}
\end{figure}

\gentab{m2/m2_data_treatment}

% ------------------------------------------------------------------------------------------------------------
\subsection{Feature-family ablation}\label{app:m2:families}

\paragraph{Design.} Each family is dropped in turn, or kept alone, in two designs
(Figure~\ref{fig:M2_03_families}; Table~\ref{tab:m2_data_treatment}, lower rows). With mirror augmentation the
family is the \R{m2/fam_families_aug_size} comparisons with all families. Augmentation zeroes the even
features, so this design tests only the mirror-sensitive part of each family. Without mirroring every feature
can carry weight; this family has \R{m2/fam_families_nomirror_size} comparisons with all families without
mirroring.

\paragraph{Results.} The \code{xcorr} family is the only one whose removal worsens the log loss after the
adjustment, both with augmentation (\Rcib{m2/lr_drop_xcorr_d_ll}, Holm $p = \R{m2/lr_drop_xcorr_p_ll_holm}$)
and without mirroring (\Rcib{m2/lrnm_drop_xcorr_d_ll}, Holm $p = \R{m2/lrnm_drop_xcorr_p_ll_holm}$). Alone, it
almost matches all families: \R{m2/lr_only_xcorr_accuracy} against \R{m2/lr_full_accuracy}, log loss
\Rcib{m2/lr_only_xcorr_d_ll} (Holm $p = \R{m2/lr_only_xcorr_p_ll_holm}$). Every other single family has a worse
log loss than all families after the adjustment, in both designs. Dropping \code{wavelet} or \code{global}
with augmentation is identical to the default, because both carry zero weight.

\paragraph{Even families and an abstaining ablation.} The \code{wavelet} family alone with augmentation has no
mirror-sensitive feature, so it returns $\pright = 0.5$ for every trial. It abstains on every trial: it has no
answered trial and no forced choice, and its \R{m2/lr_only_wavelet_accuracy} is the half credit of the
abstentions, not evidence about the family. Without mirroring the even families are tested properly: alone,
\code{wavelet} reaches \Rci{m2/lrnm_only_wavelet_accuracy} and \code{global} \Rci{m2/lrnm_only_global_accuracy},
both at chance (the CIs include one half).

\paragraph{Decision.} Keep every family (the pre-specified default). The families overlap, \code{xcorr}
carries the most information that the others do not, and the redundancy costs little in a penalised model. Its
odd correlations pair a chiral channel with a mirror-even one.

\begin{figure}[tbp]
\centering
\genfig{m2/M2_03_families}
\caption{Feature-family ablation of the L2 LR with mirror augmentation (top; numbers: mirror-sensitive / all
features of the family) and without mirroring (bottom). Left: log-loss change when one family is dropped
(paired 95\,\% CI; orange: Holm-adjusted $p < 0.05$ within the design's family; open: identical within solver
precision). Right: accuracy of each family alone; dashed: all families. Only dropping \code{xcorr} hurts; the
\code{wavelet} family alone abstains under augmentation.}
\label{fig:M2_03_families}\label{fig:app:m2-families}
\end{figure}

\FloatBarrier
% ------------------------------------------------------------------------------------------------------------
\subsection{Mirror handling}\label{app:m2:mirror}

\paragraph{Options.} (i) None. (ii) Augmentation: every training trial also enters mirrored, through the exact
feature map, with the opposite label \cite{um2017augmentation}. (iii) The antisymmetric wrapper, which forces
$\pright$ of a mirrored trial to be one minus that of the original:
\begin{equation}
  p(\mathbf{x}) = \tfrac{1}{2}\bigl(f(\mathbf{x}) + 1 - f(m(\mathbf{x}))\bigr),
  \label{eq:m2-antisym}
\end{equation}
where $f$ is the fitted model and $\mathbf{x}$ the feature vector. (iv) Both. The LR was run with all four,
LGBM with none and augmentation. The comparisons belong to the data-treatment family of
\R{m2/fam_data_treatment_size} (LR) and to a family of \R{m2/fam_mirror_lgbm_size} comparison (LGBM).

\paragraph{Evidence.} Figure~\ref{fig:M2_04_mirror}. Without mirroring, the LR changes the log loss by
\Rcib{m2/lr_mirror_none_d_ll} (Holm $p = \R{m2/lr_mirror_none_p_ll_holm}$) and the accuracy by
\Rcib{m2/lr_mirror_none_d_acc}~points (Holm $p = \R{m2/lr_mirror_none_p_acc_holm}$): not different after the
adjustment. The wrapper alone does not differ either (log loss \Rcib{m2/lr_mirror_anti_d_ll}, Holm
$p = \R{m2/lr_mirror_anti_p_ll_holm}$). Augmentation plus the wrapper is identical to augmentation alone
(largest $|\Delta\pright|$ \R{m2/lr_mirror_both_max_dp}; \R{m2/lr_mirror_both_class_changes} trials change
hand): an augmented linear model is already antisymmetric. For LGBM, dropping the augmentation worsens the log
loss by \Rcib{m2/lgbm_mirror_none_d_ll} (Holm $p = \R{m2/lgbm_mirror_none_p_ll_holm}$), with an accuracy change
of \Rcib{m2/lgbm_mirror_none_d_acc}~points (Holm $p = \R{m2/lgbm_mirror_none_p_acc_holm}$).

\paragraph{Left-handers.} Only \R{m2/mh_lr_full_left_persons} persons are left-handed, so they are shown
individually. With augmentation, the LR scores \R{m2/mh_lr_full_p11} on P11 and \R{m2/mh_lr_full_p17} on P17;
without mirroring, \R{m2/mh_lr_mirror_none_p11} and \R{m2/mh_lr_mirror_none_p17}. Two persons give no evidence
either way. The \R{m2/mh_lr_full_right_persons} right-handers score \Rci{m2/mh_lr_full_right_acc} with
augmentation and \Rci{m2/mh_lr_mirror_none_right_acc} without.

\paragraph{Decision.} Keep augmentation (pre-specified). It builds in the symmetry that the physics predicts
(Section~\ref{sec:physics}), no option beats it, and LGBM is worse without it.

\begin{figure}[tbp]
\centering
\genfig{m2/M2_04_mirror}
\caption{Mirror handling: none, augmentation, the antisymmetric wrapper (Eq.~\ref{eq:m2-antisym}) and both, for the
L2 LR (blue) and LGBM (orange; none and augmentation only). Left and middle: accuracy and log loss with 95\,\%
CIs. Right: accuracy of the \R{m2/mh_lr_full_right_persons} right-handed persons (with CI) and of each of the
\R{m2/mh_lr_full_left_persons} left-handed persons (open squares, no CI). Augmentation is kept.}
\label{fig:M2_04_mirror}\label{fig:app:m2-mirror}
\end{figure}

% ------------------------------------------------------------------------------------------------------------
\subsection{Learning curve}\label{app:m2:learning}

\paragraph{Design.} Inside every outer fold, the default LR is trained on a random subset of 4, 8, 12 or 16 of
the training persons, with \R{m2/lc_reps_planned} seeded repetitions per size; the held-out person is always
scored (Figure~\ref{fig:M2_05_learning_curve}, Table~\ref{tab:m2_learning_curve}).

\paragraph{Result.} The mean accuracy is \R{m2/lc_n04_acc_mean} (range \R{m2/lc_n04_acc_min} to
\R{m2/lc_n04_acc_max}) with 4 training persons, \R{m2/lc_n16_acc_mean} (range \R{m2/lc_n16_acc_min} to
\R{m2/lc_n16_acc_max}) with 16 and \R{m2/lc_all_acc} with all \R{m2/lc_all_persons}. The mean log loss falls
from \R{m2/lc_n04_ll_mean} to \R{m2/lc_n16_ll_mean} and \R{m2/lc_all_ll}. From 16 persons to all, the accuracy
changes by \Rb{m2/lc_last_step_acc}~points, within the range of the 16-person repetitions, while the log loss
changes by \R{m2/lc_last_step_ll}, outside their range. The accuracy has flattened while the probabilities
still sharpen: more persons would mainly improve the log loss. This reading is tentative: the full size is a
single run, and the subsets share persons.

\begin{figure}[tbp]
\centering
\genfig{m2/M2_05_learning_curve}
\caption{Learning curve of the default L2 LR: training persons subsampled inside every outer fold
(\R{m2/lc_reps_planned} seeded repetitions per size; band: range; line: mean); the held-out person is always
scored. The right-most point uses every training person (one run, 95\,\% CI). Accuracy flattens; the log loss
still falls.}
\label{fig:M2_05_learning_curve}\label{fig:app:m2-learning}
\end{figure}

\gentab{m2/m2_learning_curve}

% ------------------------------------------------------------------------------------------------------------
\subsection{What drives the predictions}\label{app:m2:importance}

\paragraph{Method.} The importance of a feature group is its share of the mean absolute contribution to the
logit on held-out persons, aggregated over the folds (Figure~\ref{fig:M2_06_importance}). For the LR the
contribution is the coefficient times the scaled feature value; for LGBM it is the tree's additive contribution
(TreeSHAP). A correlation feature is split equally between its two channels. These are associations on
held-out persons, not causal effects.

\paragraph{Result.} For the default LR, \code{xcorr} carries \R{m2/imp_lr_full_fam_xcorr}, \code{bins}
\R{m2/imp_lr_full_fam_bins}, \code{stats} \R{m2/imp_lr_full_fam_stats}, \code{fft} \R{m2/imp_lr_full_fam_fft}
and \code{cumint} \R{m2/imp_lr_full_fam_cumint}; \code{wavelet} and \code{global} carry
\R{m2/imp_lr_full_fam_wavelet} and \R{m2/imp_lr_full_fam_global}, as augmentation implies. The leading channels are the yaw rate about gravity
(\R{m2/imp_lr_full_ch_yaw_rate}), the cumulative yaw (\R{m2/imp_lr_full_ch_yaw_cum}), the device $x$
acceleration (\R{m2/imp_lr_full_ch_acc_x}) and the device $z$ rotation rate (\R{m2/imp_lr_full_ch_gyro_z}). Odd
features carry \R{m2/imp_lr_full_parity_odd}, swapped pairs \R{m2/imp_lr_full_parity_swap} and even features
\R{m2/imp_lr_full_parity_even}. For LGBM, \code{bins} (\R{m2/imp_lgbm_fam_bins}) and \code{xcorr}
(\R{m2/imp_lgbm_fam_xcorr}) lead, the cumulative yaw is the top channel (\R{m2/imp_lgbm_ch_yaw_cum}), and even
features carry \R{m2/imp_lgbm_parity_even}: trees can use them in interaction with odd ones. The yaw about
gravity leads, as the physics predicts (Section~\ref{sec:physics}; Appendix~\ref{app:m1}).

\paragraph{Caveat: the device frame.} The device-frame channels carry \R{m2/imp_lr_full_frame_device} of the
LR's contributions and \R{m2/imp_lgbm_frame_device} of LGBM's. They depend on the grip. Lane final tests this
by training on upright grips only (Appendix~\ref{app:final}): on the rotated grips the default LR falls to
\R{final/rob_m2_lr_grip} (\Rb{final/rob_m2_lr_grip_d}~points, Holm $p$~\R{final/rob_m2_lr_grip_holm}), LGBM changes by
\Rb{final/rob_m2_lgbm_grip_d}~points (Holm $p = \R{final/rob_m2_lgbm_grip_holm}$) and the stack by
\Rb{final/rob_m2_stack_grip_d}~points (Holm $p = \R{final/rob_m2_stack_grip_holm}$), each Holm-adjusted within
lane final's robustness family. Lane final therefore excludes these models as not robust.

\begin{figure}[tbp]
\centering
\genfig{m2/M2_06_importance}
\caption{What drives the predictions on held-out persons: share of the mean absolute contribution to the logit
(L2 LR: coefficient times scaled value; LGBM: TreeSHAP), by feature family (left) and by channel (right, top
10; a correlation pair is split equally between its channels). The yaw about gravity leads; the device-frame
channels carry a large share.}
\label{fig:M2_06_importance}\label{fig:app:m2-importance}
\end{figure}

\FloatBarrier
% ------------------------------------------------------------------------------------------------------------
\subsection{Calibration}\label{app:m2:calibration}

\paragraph{Method.} The expected calibration error (ECE) uses \R{m2/ece_bins} equal-mass bins. It is judged
against its perfect-calibration null band: the ECE that outcomes drawn from Bernoulli(\pright) would give with
the same predictions (trials independent, so a lower bound on the noise). No percentile bootstrap CI of the ECE
is used, because resampling inflates the ECE. The calibration slope and the CITL of a logistic recalibration
come with person-bootstrap 95\,\% CIs; perfect calibration gives 1 and 0.

\paragraph{Detection.} The nested selection has an ECE of \R{m2/cal_select_none_ece} (null 97.5th percentile
\R{m2/cal_select_none_ece_null_hi}), a slope of \Rci{m2/cal_select_none_slope} and a CITL of
\Rci{m2/cal_select_none_citl}. The default LR has an ECE of \R{m2/cal_lr_full_none_ece} (null 97.5th
percentile \R{m2/cal_lr_full_none_ece_null_hi}), a slope of \Rci{m2/cal_lr_full_none_slope} and a CITL of
\Rci{m2/cal_lr_full_none_citl}. LGBM has an ECE of \R{m2/cal_lgbm_none_ece} (null 97.5th percentile
\R{m2/cal_lgbm_none_ece_null_hi}), a slope of \Rci{m2/cal_lgbm_none_slope} and a CITL of
\Rci{m2/cal_lgbm_none_citl}; the stack an ECE of \R{m2/cal_ens_stack_none_ece} (null 97.5th percentile
\R{m2/cal_ens_stack_none_ece_null_hi}), a slope of \Rci{m2/cal_ens_stack_none_slope} and a CITL of
\Rci{m2/cal_ens_stack_none_citl} (Table~\ref{tab:m2_calibration}; Figure~\ref{fig:M2_07_calibration}). No model
shows detectable miscalibration: every slope CI includes 1, every CITL CI includes 0 and every ECE lies within
its null band.

\paragraph{Remedies.} Per model, \R{m2/cal_select_n_calibrators} calibrators were tried: Platt scaling \cite{platt1999probabilistic}, temperature
scaling \cite{guo2017calibration}, isotonic regression, beta calibration and a Platt map per platform. Each is
fitted inside every outer fold on cross-fitted predictions of the training persons and applied to the held-out
person. They are compared with the uncalibrated model by paired log-loss and Brier differences, Holm-adjusted
over the \R{m2/cal_select_n_calibrators} calibrators of a model (Figure~\ref{fig:M2_12_calibration_remedies}).
None lowers the log loss after the adjustment. The best for the selection is beta calibration
(\Rcib{m2/cal_select_beta_d_ll}, Holm $p = \R{m2/cal_select_beta_p_ll_holm}$), for the default LR beta
calibration (\Rcib{m2/cal_lr_full_beta_d_ll}, Holm $p = \R{m2/cal_lr_full_beta_p_ll_holm}$) and for LGBM
temperature scaling (\Rcib{m2/cal_lgbm_temperature_d_ll}, Holm $p = \R{m2/cal_lgbm_temperature_p_ll_holm}$).
Decision: every model stays uncalibrated.

\paragraph{Platforms.} For the selection, the Android persons (\R{m2/cal_select_none_android_n_persons}) have a
slope of \Rci{m2/cal_select_none_android_slope}, a CITL of \Rci{m2/cal_select_none_android_citl} and an ECE of
\R{m2/cal_select_none_android_ece} (null 97.5th percentile \R{m2/cal_select_none_android_ece_null_hi}). The iOS
persons (\R{m2/cal_select_none_ios_n_persons}) have a slope of \Rci{m2/cal_select_none_ios_slope}, a CITL of
\Rci{m2/cal_select_none_ios_citl} and an ECE of \R{m2/cal_select_none_ios_ece} (null 97.5th percentile
\R{m2/cal_select_none_ios_ece_null_hi}). A platform-specific Platt map against one global map changes the log
loss by \Rcib{m2/cal_select_platt_platform_vs_platt_d_ll} for the selection (Holm
$p = \R{m2/cal_select_platt_platform_vs_platt_p_ll_holm}$) and by
\Rcib{m2/cal_lr_full_platt_platform_vs_platt_d_ll} for the default LR (Holm
$p = \R{m2/cal_lr_full_platt_platform_vs_platt_p_ll_holm}$), adjusted over the
\R{m2/cal_select_platt_platform_vs_platt_family_size} analysed models (Figure~\ref{fig:M2_08_calibration_platform};
Table~\ref{tab:m2_calibration_platform}). There is no evidence that the platforms need different maps, but the
comparison is underpowered (\R{m2/cal_select_none_ios_n_persons} iOS persons).

\begin{figure}[tbp]
\centering
\genfig{m2/M2_07_calibration}
\caption{Reliability of the uncalibrated models on held-out persons (\R{m2/ece_bins} equal-mass bins; band: 95\,\% interval of
the observed share; dashed: perfect calibration). Text: ECE with the 97.5th percentile of its perfect-calibration
null, and the calibration slope with its 95\,\% CI. No model shows detectable miscalibration.}
\label{fig:M2_07_calibration}\label{fig:app:m2-calibration}
\end{figure}

\begin{figure}[tbp]
\centering
\genfig{m2/M2_12_calibration_remedies}
\caption{Post-hoc calibrators fitted inside every outer fold on cross-fitted predictions of the training
persons: paired per-person differences against the uncalibrated model in log loss (left) and Brier score
(right), 95\,\% CIs (below 0: better). Open markers: not significant after Holm over the
\R{m2/cal_select_n_calibrators} calibrators of a model, which is every case. Every model stays uncalibrated.}
\label{fig:M2_12_calibration_remedies}\label{fig:app:m2-remedies}
\end{figure}

\begin{figure}[tbp]
\centering
\genfig{m2/M2_08_calibration_platform}
\caption{Calibration per platform (Android \R{m2/cal_select_none_android_n_persons} persons, iOS
\R{m2/cal_select_none_ios_n_persons}): calibration slope with 95\,\% CI (left), ECE against its
perfect-calibration null band (middle) and the paired log-loss difference of a per-platform Platt map against
one global map (right; below 0: the per-platform map helps). No evidence for different maps; underpowered.}
\label{fig:M2_08_calibration_platform}\label{fig:app:m2-platform}
\end{figure}

\gentab{m2/m2_calibration}

\gentab{m2/m2_calibration_platform}

\FloatBarrier
% ------------------------------------------------------------------------------------------------------------
\subsection{Breakdowns and persons}\label{app:m2:breakdown}

\paragraph{Strata.} Conditions, handedness, flags and the predictor gates stratify the predictions only; they
are never inputs. A stratum gets a CI only with at least 5 persons (Figure~\ref{fig:M2_09_breakdown};
Table~\ref{tab:m2_breakdown}). The strata are descriptive, not tests of the conditions. Android
(\R{m2/bd_select_platform_android_persons} persons) reaches \Rci{m2/bd_select_platform_android_acc} and iOS
(\R{m2/bd_select_platform_ios_persons}) \Rci{m2/bd_select_platform_ios_acc}. The top-right grip reaches
\Rci{m2/bd_select_orientation_top_right_toward_acc}, against \Rci{m2/bd_select_orientation_top_left_toward_acc}
for the top-left grip and \Rci{m2/bd_select_orientation_upright_toward_acc} upright. The lowest stratum with at
least 5 persons is the twist variant: \Rci{m2/bd_select_variant_twist_acc}
(\R{m2/bd_select_variant_twist_n} trials, \R{m2/bd_select_variant_twist_persons} persons). Trials that fail the
predictor gates (\R{m2/bd_select_gate_pass_fail_n} trials of \R{m2/bd_select_gate_pass_fail_persons} persons)
reach \Rci{m2/bd_select_gate_pass_fail_acc}, against \Rci{m2/bd_select_gate_pass_pass_acc} for those that pass.

\paragraph{Small strata, person by person.} The bent-start variant has \R{m2/bd_select_variant_bent_start_persons} persons: P01
\R{m2/bdp_select_variant_bent_start_p01}, P09 \R{m2/bdp_select_variant_bent_start_p09}, P13
\R{m2/bdp_select_variant_bent_start_p13} and P15 \R{m2/bdp_select_variant_bent_start_p15}. The left-handers
score \R{m2/bdp_select_handedness_left_p11} (P11) and \R{m2/bdp_select_handedness_left_p17} (P17). The trials
flagged for a grip mismatch score \R{m2/bdp_select_grip_mismatch_flagged_p05} (P05),
\R{m2/bdp_select_grip_mismatch_flagged_p09} (P09) and \R{m2/bdp_select_grip_mismatch_flagged_p16} (P16); those
flagged for possible clipping \R{m2/bdp_select_clip_candidate_flagged_p01} (P01) and
\R{m2/bdp_select_clip_candidate_flagged_p20} (P20); the screen-rotated trials
\R{m2/bdp_select_screen_rotated_flagged_p14} (P14).

\paragraph{Persons.} Per-person accuracy, not the pooled mean, is what a new user experiences
(Figure~\ref{fig:M2_10_persons}). Under the nested selection the median person scores
\R{m2/pp_select_median}, the lowest \R{m2/pp_select_min} (\R{m2/pp_select_min_person}), and
\R{m2/select_below80} of \R{m2/pp_select_n} persons score below 80\,\%. The same person is lowest under the
default LR (\R{m2/pp_lr_full_min}, \R{m2/pp_lr_full_min_person}).

\begin{figure}[tbp]
\centering
\genfig{m2/M2_09_breakdown}
\caption{Accuracy of the nested selection by stratum (strata used after prediction only). 95\,\% CI only for
strata with at least 5 persons; smaller strata show each person (open circles; an arrow marks a person below
the axis) and the pooled value (diamond). Dashed: overall. Right: trials / persons.}
\label{fig:M2_09_breakdown}\label{fig:app:m2-breakdown}
\end{figure}

\begin{figure}[tbp]
\centering
\genfig{m2/M2_10_persons}
\caption{Accuracy of every held-out person, sorted by the nested selection (solid) with the default L2 LR
(light); colour: platform; dashed: 80\,\%. \R{m2/select_below80} of \R{m2/pp_select_n} persons fall below 80\,\%
under the nested selection.}
\label{fig:M2_10_persons}\label{fig:app:m2-persons}
\end{figure}

\gentab{m2/m2_breakdown}

\FloatBarrier
% ------------------------------------------------------------------------------------------------------------
\subsection{Several attempts and the displayed probability}\label{app:m2:attempts}

\paragraph{Question.} Production may combine k attempts of one person and hand. Two choices follow: which
aggregate decides the hand, and which probability is shown. Here k held-out primary trials of the same person
and hand stand in for k attempts. The analysis uses the \R{m2/ka_select_n_persons} persons with at least 5
trials of each hand (Figure~\ref{fig:M2_11_k_attempts}; Table~\ref{tab:m2_k_attempts}).

\paragraph{Options.} Four aggregates of the k attempts: the mean \pright, the mean logit, the sum of logits
(attempts as independent evidence) and the tempered sum $\sum_i z_i/\tauk$, where $z_i$ is the logit of attempt
$i$. The scale \tauk{} is fitted inside each outer fold on cross-fitted predictions of the training persons, so
the held-out person never informs its own \tauk.

\paragraph{Decision accuracy.} Averaging \pright{} raises the per-person balanced accuracy of the nested
selection. It is \Rci{m2/select_acc_k1} with one attempt, \Rci{m2/select_acc_k3} with three and
\Rci{m2/select_acc_k5} with five.

\paragraph{Displayed probability.} At $k = 5$ the log loss is \Rci{m2/ka_select_k5_meanp_ll} for the mean
\pright, \Rci{m2/ka_select_k5_meanlogit_ll} for the mean logit, \Rci{m2/ka_select_k5_sum_ll} for the sum of
logits and \Rci{m2/ka_select_k5_tempered_ll} for the tempered sum. The confidence gap (mean confidence minus
share correct) is \Rci{m2/ka_select_k5_meanp_gap}, \Rci{m2/ka_select_k5_meanlogit_gap},
\Rci{m2/ka_select_k5_sum_gap} and \Rci{m2/ka_select_k5_tempered_gap}, respectively: the mean \pright{} and the
mean logit are under-confident, while the CIs of the two sums include 0. Over the three pairwise comparisons at
$k = 5$ (Holm), the tempered sum has a lower log loss than the mean \pright{}
(\Rcib{m2/ka_select_k5_tempered_vs_meanp_d_ll}, Holm $p = \R{m2/ka_select_k5_tempered_vs_meanp_p_holm}$), but not
than the plain sum
(\Rcib{m2/ka_select_k5_tempered_vs_sum_d_ll}, Holm $p = \R{m2/ka_select_k5_tempered_vs_sum_p_holm}$). The
median $\tau_5$ is \R{m2/ka_select_tau_k5}: five attempts carry about \R{m2/ka_select_eff_k5} independent
attempts' worth of evidence.

\paragraph{Decision.} Decide with the mean \pright{} (the accuracy curve) and display the tempered sum, which
has the lowest log loss. The aggregates were checked by the confidence gap only, without a slope, intercept or
ECE null band, so this is weaker evidence than the single-attempt calibration above. Caveat: \tauk{} is fitted
on inner predictions from models trained on fewer persons than the final ones. If those predictions are less
confident, the tempered sum can be slightly over-confident on new persons. The production pipeline combines
attempts with its own classifier and stopping rule (Appendix~\ref{app:prod}).

\begin{figure}[tbp]
\centering
\genfig{m2/M2_11_k_attempts}
\caption{k attempts of one person and hand (\R{m2/ka_select_n_persons} persons). Left: per-person balanced
accuracy of the mean \pright{} (band: 95\,\% CI of the nested selection). Middle and right: log loss and
confidence gap (above 0: over-confident) of four aggregates for the nested selection, with 95\,\% CIs. Decide
with the mean \pright; display the tempered sum.}
\label{fig:M2_11_k_attempts}\label{fig:app:m2-attempts}
\end{figure}

\gentab{m2/m2_k_attempts}

% ------------------------------------------------------------------------------------------------------------
\subsection{Caveats}\label{app:m2:caveats}

\begin{itemize}
  \item With \R{describe/n_persons} persons every CI is wide; differences of a point or two between pipelines are
  not resolved.
  \item Only the nested selection and the pre-specified default are headline-grade. The model comparison, the
  calibrators and the ablations are exploratory; every result records its number of configurations compared.
  \item The features use the oracle (app) segmentation. Production sees detected windows, and the approach
  features inherit any segmentation error there.
  \item The free-choice trials use one hand only: they give a right-hand recall, not a calibration or accuracy
  check over both hands.
  \item The experiments use \R{m2/n_boot} person resamples, the derived intervals (strata, platforms, slopes)
  \R{m2/n_boot_derived} and the multi-attempt intervals \R{m2/n_boot_attempts}.
  \item Open: does per-trial amplitude normalisation hold on detected windows?
\end{itemize}

\FloatBarrier
